\subsection{初始一致结构}

命题 4 设 X 是一个集合，\((\mathrm{Y}_{t})_{t\in I}\) 是一族一致空间，且对每个 \(t\in I\)，\(f_{t}\) 是 X 到 \(Y_{t}\) 的映射。对每个 \(t\in I\)，令 \(g_{t}\) 表示 \(f_{t}\times f_{t}\)。设 S 为 \(X\times X\) 中形如 \(\overline{g}_{t}(V_{t})\) 的子集组成的集，其中 \(t\in I\) 而 \(V_{t}\) 是 \(Y_{t}\) 的一致邻域；又设 B 为一切有限交

\[
\mathbf {U} \left(\mathrm{V} _ {\iota_ {1}}, \dots , \mathrm{V} _ {\iota_ {n}}\right) = _ {i} ^ {- 1} g _ {i} \left(\mathrm{V} _ {\iota_ {1}}\right) \cap \dots \cap^ {- 1} g _ {\iota_ {n}} \left(\mathrm{V} _ {\iota_ {n}}\right)\tag{1}
\]

组成的集，其中每个成员由 \(\mathfrak{S}\) 中的集合构成。于是 \(\mathfrak{B}\) 是某个一致结构 \(\mathfrak{U}\)（\(\mathbf{X}\) 上的）的一致邻域基本系，该一致结构是 \(\mathbf{X}\) 关于族 \((f_{\mathfrak{i}})\) 的初始一致结构（《集合论》，第四章，§ 2，第 3 目），特别地，\(\mathfrak{U}\) 是 \(\mathbf{X}\) 上使所有映射 \(f_{\mathfrak{i}}\) 都一致连续的最粗一致结构。进而言之，设 \(h\) 是一致空间 \(\mathbf{Z}\) 到 \(\mathbf{X}\) 的映射；则 \(h\) 一致连续（当 \(\mathbf{X}\) 赋有一致结构 \(\mathfrak{U}\) 时），当且仅当每个映射 \(f_{\mathfrak{i}} \circ h\) 都一致连续。

立即可以看出 \(\mathfrak{B}\) 满足公理 \((\mathbf{B}_1)\mathbf{Y}\) 与 \((\mathbf{U}_1')\)。若 \(\mathbf{W}_t = \overline{g}_t^1 (\mathbf{V}_t)\)，则 \(\overline{\mathbf{W}}_t = \overline{g}_t^1 (\overline{\mathbf{V}}_t^1)\) 且 \(\dot{\mathbf{W}}_t = \overline{g}_t^1 (\overline{\mathbf{V}}_t^2)\)；于是 \(\mathfrak{B}\) 也满足公理 \((\mathbf{U}_{\mathrm{II}}^t)\) 与 \((\mathbf{U}_{\mathrm{III}}^t)\)，从而是某个一致结构 \(\mathcal{U}\)（\(\mathbf{X}\) 上的）的一致邻域基本系。此外，由 \(\mathcal{U}\) 的定义以及第 1 目的定义 1 立即可知，\(f_{t}\) 对每个指标 \(i\in I\) 一致连续；因此（第 1 目命题 2）\(f_{t}\circ h\) 对每个 \(i\in I\) 都一致连续，只要 \(h\) 一致连续。反之，设 \(f_{t}\circ h\) 对每个 \(i\in I\) 都一致连续，并考虑一个集合 \(\mathrm{U}(\mathrm{V}_{i_t},\dots,\mathrm{V}_{i_n})\)；由假设，对每个 \(k\)（\(1\leqslant k\leqslant n\)），存在 \(\mathbf{W}_k\)（\(\mathbf{Z}\) 的一致邻域），使关系 \((z,z^{\prime})\in \mathbf{W}_k\) 蕴含 \([f_{i_k}(h(z)),f_{i_k}(h(z')))]\in \mathbf{V}_k\)；若

\[
\mathbf {W} = \bigcap_ {k} \mathbf {W} _ {k},
\]

则这 \(n\) 个关系只要 \(z\) 与 \(z'\) 是 W-邻近的便同时满足，从而有 \((h(z), h(z')) \in \mathrm{U}(\mathrm{V}_{\iota_1}, \ldots, \mathrm{V}_{\iota_n})\)，证明完成。

推论 使诸 \(f_{i}\) 一致连续的最粗一致结构 U 在 X 上诱导的拓扑，也是使诸 \(f_{i}\) 连续的最粗拓扑。

这是后一拓扑中一点的邻域定义（第一章 § 2 第 3 目命题 4）的直接推论。

初始结构的一般性质（《集合论》，第四章，§ 2，第 3 目，判别准则 CST 10）特别蕴含下列传递性质：

命题 5 设 X 是一个集合，\((Z_{\iota})_{\iota \in I}\) 是一族一致空间，\((J_{\lambda})_{\lambda \in L}\) 是 I 的一个分拆，\((Y_{\lambda})_{\lambda \in L}\) 是以 L 为指标集的一族集合。对每个 \(\lambda \in L\)，设 \(h_{\lambda}\) 是 X 到 \(Y_{\lambda}\) 的映射；对每个 \(\lambda \in L\) 与每个 \(\iota \in J_{\lambda}\)，设 \(g_{\iota \lambda}\) 是 \(Y_{\lambda}\) 到 \(Z_{\iota}\) 的映射；并令 \(f_{\iota} = g_{\iota \lambda} \circ h_{\lambda}\)。设每个 \(Y_{\lambda}\) 赋有使诸映射 \(g_{\iota \lambda} (\iota \in J_{\lambda})\) 一致连续的最粗一致结构。则使诸 \(f_{\iota}\) 一致连续的 X 上最粗的一致结构，与使诸 \(h_{\lambda}\) 一致连续的 X 上最粗的一致结构是同一个。

